% Regression test: Leslie Lamport's pf2, the structured-proof package that
% issue #4 was reported against. It is not on CTAN (it is distributed from
% Lamport's own site), so this fixture is skipped unless pf2.sty is findable;
% see the note in examples/Makefile.
%
% pf2 writes a proof as a tree: every step may be proved by a nested proof
% environment, arbitrarily deep, and a reference made at any depth is a
% dependency of the one result the whole tree establishes. That is the shape
% \proofof has to reach through, and the shape it did not: a nested proof took
% the result preceding the proof around it, so the references below the top
% level went to whichever result happened to be declared last.
%
% pf2 defines its own \newenvironment{proof}, which collides with the proof of
% amsthm that proofgraph loads. The document therefore does what the reporter
% did, renaming pf2's environment -- here by lending it a second name rather
% than by editing pf2.sty -- and declares the new name with \proofgraphproofenv.
% amsthm's proof is put back afterwards and used too, so both kinds of proof
% live in one document and the graph must be right for each.
%
% Single-pass and bibtex-free for determinism.
\documentclass{article}
\usepackage{amsthm}
\usepackage{proofgraph}

% Free the names pf2 insists on defining (proof, from amsthm, and qed), keeping
% amsthm's to put back after. No @ in these names: we are in the preamble of a
% document, where @ is not a letter.
\let\amsproof\proof
\let\amsendproof\endproof
\let\amsqed\qed
\let\proof\relax
\let\endproof\relax
\let\qed\relax
\usepackage{pf2}
% pf2's proof, under a name of its own, and amsthm's proof and qed back in
% place. The structured proofs below use \step, not pf2's \qed.
\let\structuredproof\proof
\let\endstructuredproof\endproof
\let\proof\amsproof
\let\endproof\amsendproof
\let\qed\amsqed

\proofgraphproofenv{structuredproof}

\newtheorem{theorem}{Theorem}
\newtheorem{lemma}[theorem]{Lemma}

\begin{document}

\begin{lemma}\label{lem:a}A.\end{lemma}
\begin{lemma}\label{lem:b}B.\end{lemma}
\begin{lemma}\label{lem:c}C.\end{lemma}
\begin{lemma}\label{lem:d}D.\end{lemma}
\begin{theorem}\label{thm:tree}Proved by a structured proof.\end{theorem}
\begin{theorem}\label{thm:plain}Proved by an amsthm proof.\end{theorem}
\begin{theorem}\label{thm:last}Declared last, proves nothing.\end{theorem}

Here is the proof of Theorem~\ref{thm:tree}:

% A pf2 proof three levels deep, detached from its result and pinned by
% \proofof: every reference in the tree belongs to thm:tree.
\begin{structuredproof}
  \proofof{thm:tree}
  \step{s1}{By Lemma~\ref{lem:a}.}
  \begin{structuredproof}
    \step{s11}{By Lemma~\ref{lem:b}.}
    \begin{structuredproof}
      \step{s111}{By Lemma~\ref{lem:c}.}
    \end{structuredproof}
  \end{structuredproof}
\end{structuredproof}

% amsthm's proof still works, and still attaches to the result it follows.
\begin{proof}[Proof of Theorem~\ref{thm:plain}]
  By Lemma~\ref{lem:d}.
\end{proof}

\end{document}
